% GENERATED FILE. Edit canonical lessons, then run npm run generate:books.
% canonical-source-hash: fnv1a64:cc563d665732af74
% canonical-lessons: LA-C352-mola, LA-C352-cribrum, LA-C352-subigo, LA-C352-fermentum, LA-C352-clibanus

\chapter{Pānis domesticus}
\label{ch:la-panis-domesticus}
\hlchaptermodality{352}

\begin{chapteropening}
\textbf{By the end of this chapter:} \emph{I can tell how bread was made: the millstone, the sieve, to knead, the yeast and the baking pot.}

Complete the last lesson of chapter 352: I can tell how bread was made: the millstone, the sieve, to knead, the yeast and the baking pot.
\end{chapteropening}

\section[mola, molae]{mola, molae — a millstone; a mill; salted meal}

\label{lesson:LA-C352-mola}



\begin{quote}
Before the new one: say the Latin for a pestle, then the Latin for to grind.
\end{quote}

\subsection*{You'll want to know: mola, molae}
\textbf{mola, molae} — \textquotedblleft{}a millstone; a mill; salted meal\textquotedblright{}.

\begin{quote}
Asinus molam in pistrīnō vertēbat.
\end{quote}

\begin{quote}
\emph{A donkey was turning the millstone in the bakery.}
\end{quote}

Feminine; the plural \emph{molae} is a mill. \emph{Mola salsa}, salted meal, was sprinkled on victims. English \emph{molar}, the grinding tooth.

\begin{cousinweb}
\textbf{mola} is inherited from the Proto-Indo-European root \emph{*melh\textsubscript{2}-}, \textquotedblleft{}to grind\textquotedblright{} (EtymDB). English \emph{molar} comes from it, and English \emph{meal}, \textquotedblleft{}flour\textquotedblright{}, is a cousin.
\end{cousinweb}

\begin{grammarlens}[title={What you've built}]
One more word for telling, explaining and describing at B1.
\end{grammarlens}

\subsection*{Your turn}
\begin{itemize}
\raggedright
  \item \emph{Say it:} \emph{mola, molae}
  \item \emph{Say it:} \emph{mola, molae}, once more
  \item \emph{Say it:} a sentence of your own with \emph{mola, molae}
  \item \emph{Recall:} say the Latin for a pestle, then the Latin for to grind, then say \emph{mola, molae} again
\end{itemize}

\begin{tcolorbox}[breakable,colback=teal!4,colframe=teal!35!black,title={Before you move on}]
What does mola, molae mean? (\textquotedblleft{}A millstone; a mill; salted meal\textquotedblright{}.) Say it once more.
\end{tcolorbox}
\section[crībrum, crībrī]{crībrum, crībrī — a sieve}

\label{lesson:LA-C352-cribrum}



\begin{quote}
Before the new one: say the Latin for to grind, then the Latin for a millstone.
\end{quote}

\subsection*{You'll want to know: crībrum, crībrī}
\textbf{crībrum, crībrī} — \textquotedblleft{}a sieve\textquotedblright{}.

\begin{quote}
Pistor farīnam crībrō pūrgābat.
\end{quote}

\begin{quote}
\emph{The baker was cleaning the flour with a sieve.}
\end{quote}

Neuter. Pouring water into a sieve was a Roman picture of a hopeless task.

\begin{cousinweb}
\textbf{crībrum} is inherited from a Proto-Indo-European root meaning \textquotedblleft{}to sift\textquotedblright{} (EtymDB), also seen, says Lewis \& Short, in \emph{cernō}, \textquotedblleft{}to sift, to separate\textquotedblright{}.
\end{cousinweb}

\begin{grammarlens}[title={What you've built}]
One more word for telling, explaining and describing at B1.
\end{grammarlens}

\subsection*{Your turn}
\begin{itemize}
\raggedright
  \item \emph{Say it:} \emph{crībrum, crībrī}
  \item \emph{Say it:} \emph{crībrum, crībrī}, once more
  \item \emph{Say it:} a sentence of your own with \emph{crībrum, crībrī}
  \item \emph{Recall:} say the Latin for to grind, then the Latin for a millstone, then say \emph{crībrum, crībrī} again
\end{itemize}

\begin{tcolorbox}[breakable,colback=teal!4,colframe=teal!35!black,title={Before you move on}]
What does crībrum, crībrī mean? (\textquotedblleft{}A sieve\textquotedblright{}.) Say it once more.
\end{tcolorbox}
\section[subigō, subigere]{subigō, subigere — to knead; to subdue}

\label{lesson:LA-C352-subigo}



\begin{quote}
Before the new one: say the Latin for a millstone, then the Latin for a sieve.
\end{quote}

\subsection*{You'll want to know: subigō, subigere}
\textbf{subigō, subigere} — \textquotedblleft{}to knead; to subdue\textquotedblright{}.

\begin{quote}
Pistor farīnam et aquam subēgit.
\end{quote}

\begin{quote}
\emph{The baker kneaded the flour with water.}
\end{quote}

\emph{subigō, subigere, subēgī, subāctum}: \emph{sub} + \emph{agō}. Also to bring under one's power.

\begin{cousinweb}
\textbf{subigō} is built on \emph{agō}, \textquotedblleft{}to drive\textquotedblright{}, with \emph{sub-}, \textquotedblleft{}under\textquotedblright{} (EtymDB, Lewis \& Short): to work something under, to knead, and to subdue.
\end{cousinweb}

\begin{grammarlens}[title={What you've built}]
One more word for telling, explaining and describing at B1.
\end{grammarlens}

\subsection*{Your turn}
\begin{itemize}
\raggedright
  \item \emph{Say it:} \emph{subigō, subigere}
  \item \emph{Say it:} \emph{subigō, subigere}, once more
  \item \emph{Say it:} a sentence of your own with \emph{subigō, subigere}
  \item \emph{Recall:} say the Latin for a millstone, then the Latin for a sieve, then say \emph{subigō, subigere} again
\end{itemize}

\begin{tcolorbox}[breakable,colback=teal!4,colframe=teal!35!black,title={Before you move on}]
What does subigō, subigere mean? (\textquotedblleft{}To knead; to subdue\textquotedblright{}.) Say it once more.
\end{tcolorbox}
\section[fermentum, fermentī]{fermentum, fermentī — leaven, yeast}

\label{lesson:LA-C352-fermentum}



\begin{quote}
Before the new one: say the Latin for a sieve, then the Latin for to knead.
\end{quote}

\subsection*{You'll want to know: fermentum, fermentī}
\textbf{fermentum, fermentī} — \textquotedblleft{}leaven, yeast\textquotedblright{}.

\begin{quote}
Sine fermentō pānis dūrus fit.
\end{quote}

\begin{quote}
\emph{Without leaven bread turns out hard.}
\end{quote}

Neuter, from a verb meaning 'to boil, to bubble'. English \emph{ferment}.

\begin{cousinweb}
\textbf{fermentum} is built on \emph{ferveō}, \textquotedblleft{}to boil, to bubble\textquotedblright{} (EtymDB): what makes dough bubble. English \emph{ferment} comes from it.
\end{cousinweb}

\begin{grammarlens}[title={What you've built}]
One more word for telling, explaining and describing at B1.
\end{grammarlens}

\subsection*{Your turn}
\begin{itemize}
\raggedright
  \item \emph{Say it:} \emph{fermentum, fermentī}
  \item \emph{Say it:} \emph{fermentum, fermentī}, once more
  \item \emph{Say it:} a sentence of your own with \emph{fermentum, fermentī}
  \item \emph{Recall:} say the Latin for a sieve, then the Latin for to knead, then say \emph{fermentum, fermentī} again
\end{itemize}

\begin{tcolorbox}[breakable,colback=teal!4,colframe=teal!35!black,title={Before you move on}]
What does fermentum, fermentī mean? (\textquotedblleft{}Leaven, yeast\textquotedblright{}.) Say it once more.
\end{tcolorbox}
\section[clībanus, clībanī]{clībanus, clībanī — a baking pot, a small oven}

\label{lesson:LA-C352-clibanus}



\begin{quote}
Before the new one: say the Latin for to knead, then the Latin for leaven.
\end{quote}

\subsection*{You'll want to know: clībanus, clībanī}
\textbf{clībanus, clībanī} — \textquotedblleft{}a baking pot, a small oven\textquotedblright{}.

\begin{quote}
Pānem in clībanō coxit.
\end{quote}

\begin{quote}
\emph{He baked the bread in a baking pot.}
\end{quote}

Masculine, a Greek loan: a clay or iron pot heated with coals, wider at the bottom.

\begin{cousinweb}
\textbf{clībanus} was borrowed from Greek \emph{klibanos}, \textquotedblleft{}a baker's oven, a covered pan\textquotedblright{} (EtymDB).
\end{cousinweb}

\begin{grammarlens}[title={What you've built}]
One more word for telling, explaining and describing at B1.
\end{grammarlens}

\subsection*{Your turn}
\begin{itemize}
\raggedright
  \item \emph{Say it:} \emph{clībanus, clībanī}
  \item \emph{Say it:} \emph{clībanus, clībanī}, once more
  \item \emph{Say it:} a sentence of your own with \emph{clībanus, clībanī}
  \item \emph{Recall:} say the Latin for to knead, then the Latin for leaven, then say \emph{clībanus, clībanī} again
\end{itemize}

\begin{tcolorbox}[breakable,colback=teal!4,colframe=teal!35!black,title={Before you move on}]
What does clībanus, clībanī mean? (\textquotedblleft{}A baking pot, a small oven\textquotedblright{}.) Say it once more.
\end{tcolorbox}
